\documentclass[12pt, leqno]{article}
\usepackage[english,russian]{babel}
\usepackage{fontspec}
\usepackage[utf8]{inputenc}
\usepackage{epigraph}
\usepackage{tabularx}
\usepackage[leqno]{amsmath}
\setmainfont{Times New Roman}
\setlength\parindent{0pt}
\usepackage[colorlinks=true,urlcolor=blue]{hyperref}        
\usepackage{minted} 
\usepackage[xetex,final]{graphicx}
\usepackage{float}%"Плавающие" картинки
\usepackage{wrapfig}%Обтекание фигур (таблиц, картинок и прочег
\usepackage{epstopdf}
\usepackage{amssymb}
 
\begin{document}

\section*{Числа Каталана. Рекурсия}
\begin{flushleft}
\textbf{Вопросы рекурсии:} Рекурентные формулы и реализация
\end{flushleft}

Числа Каталана удовлетворяют рекуррентному соотношению:

$ C_{0}=1\quad  $ и $\quad C_{n}=\sum _{i=0}^{n-1}C_{i}C_{n-1-i}}$ для $ n\geqslant 1 $

Также существует более простое рекуррентное соотношение (которое мы сдесь и используем): 

$ C_{0}=1\quad  $ и $\quad C_{n}={\frac {2(2n-1)}{n+1}}C_{n-1}$


 
\clearpage 
 
Дальше можно смотреть по ссылке:

\href{https://ru.wikipedia.org/wiki/Числа_Каталана}{Числа_Каталана}

\noindent Образец кода на лисп (sbcl):  
\begin{minted}[fontsize=\small]{lisp}

;;;;
;;;; Compute Catalan Number with recursion (вычисление чисел Каталана рекурентным способом)
;;;; https://ru.wikipedia.org/wiki/Числа_Каталана
;;;; http://www.math.bas.bg/bantchev/misc/catalan.pdf
;;;; https://informatics.msk.ru/mod/book/view.php?id=266&chapterid=58
;;;; https://oeis.org/A000108
;;;; https://www.jdoodle.com/execute-clisp-online
;;;; https://www.tutorialspoint.com/lisp/lisp_constants.htm
;;;; 

;; C_{0}=1\quad
;; \quad C_{n}=\sum _{i=0}^{n-1}C_{i}C_{n-1-i}

(defun CatalansNumbers(i)
    (cond ((= i 0) 1)
          ((= i 1) 1)
          (t (* (/ (+ (* 4 (- i 1)) 2) ( + (- i 1) 2))  (CatalansNumbers  (- i 1))))
    )
)

(defun domain(n acc)
    (cond ((< n 0) acc)
          (t (domain (- n 1) (cons n acc))))
)

(defun  computeFromDomain(n)
    (mapcar (lambda (x) (CatalansNumbers x)) (domain n nil))
)

(defun cli/parameter()
    (let (
        (args sb-ext:*posix-argv*))
        (car (cdr args))))

;; main print
;; input   ./exec6.lisp  <n = 6>
(defun main()
 (terpri)
 (format T "~a" (computeFromDomain (parse-integer (cli/parameter))))
)

(main)

\end{minted}  

Ниже представлен результат:
 
1, 1, 2, 5, 14, 42, 132, 429, 1430, 4862, 16796, 58786, 208012, 742900, 2674440, 9694845, 35357670, 129644790, 477638700, 1767263190, 6564120420, 24466267020, 91482563640, 343059613650, 1289904147324, 4861946401452 

 

  

\end{document}